% SHARD-ID: AQM-79-QSA-GEOMETRIC-FIELD-HOM-X-V
% SHARD-TITLE: Geometric Field Association Hom(X,V)
% SHARD-SUMMARY: Reviews the geometric field-space branch of the QSA catalogue without asserting a general Hom(X,V) quantisation.
% SHARD-SUMMARY: Separates all-map, finite-support, sheaf-section, regular-function, and rational-truncation choices using only local anchors.
% SHARD-SUMMARY: Records how the structure of X constrains field spaces before the finite Weyl-Heisenberg boundary of AQM-78.
% SHARD-KEYWORDS: quantum-system association, Hom(X,V), Map(X,V), sheaf sections, regular functions, finite support, observable nets

\section{Geometric Field Association \texorpdfstring{\(\operatorname{Hom}(X,V)\)}{Hom(X,V)}}
\label{sec:qsa-geometric-field-hom-x-v}

\subsection{Status and Local Anchors}

\textbf{Status: Review.}  This is QSA Step 15 from
\path{docs/quantum_system_associations/PLAN.md}.  It adds no source,
convention entry, run artifact, or general functorial quantisation theorem.
Its purpose is to place the existing geometric field-space shards into the
QSA catalogue.

The title uses \(\operatorname{Hom}(X,V)\) as the prompt label for this
branch.  In the local conventions, this is deliberately not the default
object.  By \texttt{CONVENTIONS.md} convention \texttt{(q)}, a finite set
uses \(\operatorname{Map}(X,V)\); the word \(\operatorname{Hom}\) is reserved
for a specified structure-preserving class of maps.  The QSA vocabulary and
non-functoriality rule are AQM-65 and convention \texttt{(ap)}.  The
point-choice boundary is AQM-75, and the finite phase-space representation
boundary is AQM-78.

The field-space anchors reviewed here are AQM-14--AQM-18, AQM-24--AQM-30,
and AQM-49--AQM-52.  The relevant conventions are \texttt{(q)} for finite
map spaces and radicals, \texttt{(r)} for projective sheaf fields,
\texttt{(v)}--\texttt{(w)} for open-local observable nets and their
presheaf/cosheaf distinction, \texttt{(y)} for closed-point finite-residue
supports, and \texttt{(ak)}--\texttt{(al)} for regular-function profiles,
reduced smearings, and affine action.

\subsection{Choice Table}

Every row below is constrained by a convention or prior shard; no row is a universal \(\operatorname{Hom}(X,V)\) rule.
\begin{center}
\small\setlength{\tabcolsep}{3pt}
\begin{tabular}{p{0.18\linewidth}p{0.22\linewidth}p{0.19\linewidth}p{0.19\linewidth}p{0.12\linewidth}}
\toprule
Field-space choice & How \(X\) constrains it & Operator/Weyl status &
Radical or support condition & Anchor \\
\midrule
\texttt{Review:} all maps on a finite set &
Only the finite point labels of \(X\) are used; extra structure is ignored
unless a subspace is chosen. &
\texttt{Checked} finite Weyl labels after the pointwise form is reduced. &
Use \(E\leq\operatorname{Map}(X,V)\) and quotient by
\(\operatorname{rad}(E)\). &
AQM-14, AQM-71, \texttt{(q)} \\
\addlinespace
\texttt{Review:} structured finite subspaces &
Linear, affine, or polynomial structure on \(X\) may choose a smaller
evaluated subspace \(E\leq\operatorname{Map}(X,V)\). &
\texttt{Checked} only for the local \(\mathbb F_3\) rows already run. &
The same radical test applies; evaluation rank is not the same as symplectic
rank. &
AQM-14, arithmetic-field run bundle \\
\addlinespace
\texttt{Review:} finite closed supports &
For affine finite-type \(\mathbb F_q\)-schemes, only closed points are finite
residue sites; non-closed points still control topology. &
\texttt{Checked/Review} quasi-local Weyl net over finite supports. &
Labels have finite closed support:
\(\bigoplus_{x\in U\cap |X|_{\rm cl}}\kappa(x)^2\). &
AQM-28--AQM-30, \texttt{(y)} \\
\addlinespace
\texttt{Review:} rational truncations &
A selected finite set \(X(k)\) or finite rational support is used instead of
all closed points. &
\texttt{Checked/Review} finite subtheory, not the full scheme theory. &
Finite by choice of support; non-rational closed points are omitted. &
AQM-16, AQM-30, AQM-75 \\
\addlinespace
\texttt{Review:} projective sheaf sections &
Projective geometry makes regular maps to affine targets too small, so a
sheaf \(\mathcal L\) is chosen. &
\texttt{Checked} for \(\mathbf P^1_{\mathbb F_3}\) and
\(\mathcal O(d)\), \(0\leq d\leq4\). &
\(E=H^0(X,\mathcal L)\otimes V\), evaluated on chosen finite points, then
quotiented by the radical. &
AQM-15--AQM-18, \texttt{(r)} \\
\addlinespace
\texttt{Review:} regular-function profiles &
For \(X=\Spec\mathbb F_q[x]\), geometry constrains functions to
\(\Gamma(D(h),\mathcal O_X)=\mathbb F_q[x]_h\). &
\texttt{Gap} as a full quasi-local operator except for zero profile;
\texttt{Checked} for finite reduced smearings. &
Nonzero profiles have cofinite support; a polynomial \(h\) gives a finite
support \(Z_{\rm cl}(h)\) for \(W_h^{\rm red}\). &
AQM-49--AQM-52, \texttt{(ak)}--\texttt{(al)} \\
\addlinespace
\texttt{Review:} observable nets &
Zariski opens and closed supports control which finite labels are allowed in
each local algebra. &
\texttt{Checked/Review} open-local isotone net; not an ordinary cosheaf of
unital \(*\)-algebras. &
Finite-support labels glue additively; observable algebras glue by generated
commuting tensor-local subalgebras. &
AQM-24--AQM-27, \texttt{(v)}--\texttt{(w)} \\
\bottomrule
\end{tabular}
\end{center}

\subsection{Finite Maps and Closed-Point Fields}

\textbf{Review.}  The finite-set model of AQM-14 starts with a finite point set
\(X\), a finite symplectic target \((V,\omega_V)\), and a chosen
\(\mathbb F_q\)-linear subspace
\[
  E\leq \operatorname{Map}(X,V).
\]
This is not an unsourced \(\operatorname{Hom}(X,V)\).  AQM-14 D1--D4 and
convention \texttt{(q)} fix the pointwise form
\[
  \Omega_E(\phi,\psi)=
  \sum_{x\in X}w_x\,\omega_V(\phi(x),\psi(x)),
\]
with all weights equal to \(1\) in the checked runs.  AQM-14 L4--L5 and
AQM-71 record the required reduction
\[
  E_{\rm red}=E/\operatorname{rad}(E).
\]
Thus the structure of \(X\) constrains the field space only when the shard
has already named a structure-sensitive subspace, such as the linear, affine,
or bounded-degree polynomial examples evaluated in AQM-14.  The finite Weyl
step applies to \(E_{\rm red}\), not to an arbitrary unreduced display space.

\textbf{Review.}  The closed-point affine branch changes the point choice.
By AQM-28 and convention \texttt{(y)}, if
\[
  X=\Spec R
\]
is affine of finite type over \(\mathbb F_q\), the finite Weyl sites are the
closed points \(x\in |X|_{\rm cl}\).  Each such point has finite residue field
\(\kappa(x)\), local label space \(\kappa(x)^2\), carrier
\(\ell^2(\kappa(x))\), and local matrix algebra.  On an open \(U\),
\[
  E(U)=\bigoplus_{x\in U\cap |X|_{\rm cl}}\kappa(x)^2
\]
means finite closed support.  AQM-29 applies this to \(\Spec\mathbb F_q[t]\),
where each monic irreducible of degree \(d\) gives one \(q^d\)-level qudit.
AQM-30 applies it to \(\Spec\mathbb F_q[x,y]\), where rational points give
only a finite \(q^2\)-site truncation of the full closed-point theory.

\textbf{Gap.}  Generic and curve-generic points are not finite Weyl sites in
this branch.  AQM-29 and AQM-30 keep them as topological and geometric points,
but no finite-residue field operator is attached to them.  AQM-75 records
this as a point-selection boundary, and AQM-78 does not remove it: AQM-78
starts only after a finite phase-label object has already been fixed.

\subsection{Sheaf Sections and Regular Functions}

\textbf{Review.}  AQM-15 and convention \texttt{(r)} record why a projective
geometric field space is not obtained by silently reading
\(\operatorname{Hom}(X,V)\) as regular maps to an affine vector space.  For a
proper geometrically integral projective variety, the local Stacks locator in
AQM-15 gives only constant global regular functions.  The projective
replacement used in the report is instead
\[
  E(X,\mathcal L,V)=H^0(X,\mathcal L)\otimes_{\mathbb F_q}V,
\]
with the sheaf \(\mathcal L\) chosen as part of the datum.

AQM-16 checks this for
\[
  X=\mathbf P^1_{\mathbb F_3},\qquad
  \mathcal L=\mathcal O(d),\qquad V=\mathbb F_3^2,
\]
using the rational-point order and evaluation convention in
\texttt{(r)}.  The \(d=1\) row has independent evaluation rows but total
radical for the all-weights-one scalar pairing, while \(d=4\) has a
one-dimensional scalar evaluation kernel and a two-dimensional field radical.
These are checked examples of the same rule: sections first, chosen
evaluations second, radical quotient third, Weyl labels fourth.

\textbf{Review.}  AQM-17 and AQM-18 refine the finite point value: the stalk
is the local field germ, while the fiber or residue value is what the finite
Weyl pairing samples.  For \(\mathcal O(d)\otimes\mathbb F_3^2\), the stalk
reduces to \(V=\mathbb F_3^2\) only after a chosen local frame is reduced
modulo the maximal ideal.

\textbf{Review.}  The affine-line regular-function branch in AQM-49--AQM-52
is another structure-sensitive replacement for a vague Hom-space.  On
\(U=D(h)\subset\Spec\mathbb F_q[x]\), a pair
\((F,G)\in\mathbb F_q[x]_h^2\) gives a residue profile
\[
  \operatorname{ev}_U(F,G)_\pi=(F\bmod\pi,\ G\bmod\pi)\in\kappa(x_\pi)^2.
\]
By AQM-49 and convention \texttt{(ak)}, a nonzero profile on a nonempty
standard open has cofinite closed support, so it is generally a formal
infinite-volume profile rather than a quasi-local Weyl operator.

\textbf{Checked.}  A nonzero polynomial \(h\) has a different role: it
selects the finite closed support
\[
  Z_{\rm cl}(h)=\{x_\pi:\pi\mid h\}.
\]
Then AQM-49 defines the reduced smeared label and operator
\[
  W_h^{\rm red}(F,G)
  =
  W_{Z_{\rm cl}(h)}
  \bigl((F\bmod\pi,\ G\bmod\pi)_{\pi\mid h}\bigr).
\]
AQM-50 checks this line by line over \(\mathbb F_3\), including rational
two-qutrit supports, a quadratic \(9\)-level site, and the repeated-factor
case where multiplicities are invisible to the reduced support net.
AQM-51--AQM-52 and convention \texttt{(al)} check that affine symmetries move
regular functions, finite supports, and reduced smearings compatibly, while
preserving the profile/operator distinction.

\subsection{Observable Nets and the AQM-78 Boundary}

\textbf{Review.}  The observable-net layer is not a new restriction of
\(\operatorname{Hom}(X,V)\); it records where finite labels are supported.
AQM-24 and convention \texttt{(v)} use the open-local assignment
\[
  U\longmapsto\mathcal A(U),
\]
with \(U\subset V\) giving tensor-by-identity inclusions.  The
closed-complement assignment is valid only as a named closed-support layer,
because its arrows are contravariant in opens.

AQM-25 builds the algebraic quasi-local algebra as the filtered colimit over
quasi-compact opens.  AQM-26 and convention \texttt{(w)} separate the
cosheaf-like label layer
\[
  E(U)=\bigoplus_{x\in U}E_x
\]
from the observable layer, an isotone local net of unital \(*\)-algebras.
AQM-27 records the state-side reversal as a presheaf of states.  These net
statements support finite-support observables; they do not create
infinite-volume Weyl operators for regular-function profiles.

\textbf{Review.}  AQM-78 begins after one of the rows above has produced a
finite phase label with commutator data.  It supplies the Weyl-Heisenberg
projective representation and observable algebra for finite phase spaces.
Thus Step 15 does not claim that every geometric
\(\operatorname{Hom}(X,V)\) has been quantized.  It says only that the report
already contains several locally constrained field-space choices, and that
AQM-78 is the common finite Weyl boundary once the label object is finite and
the radical/support issues have been resolved.

\subsection{Next-Step Gaps}

\begin{center}
\small\setlength{\tabcolsep}{3pt}
\begin{tabular}{p{0.16\linewidth}p{0.25\linewidth}p{0.37\linewidth}p{0.12\linewidth}}
\toprule
Next step & What this shard supplies & Remaining condition & Status \\
\midrule
QSA Step 16: finite-ring example matrix &
Rows can cite whether a field-space entry is all-map, finite-support
closed-point, regular-function profile, rational truncation, or sheaf-section
style. &
Each finite-ring row still needs its own point choice, residue/cotangent or
function-space convention, and local evidence; no populated finite-ring
database claim is imported here. &
\texttt{Gap} \\
\addlinespace
QSA Step 17: residue-field site association &
AQM-23 and AQM-28--AQM-30 already show variable finite residue fields and
finite-support direct sums. &
A finite-ring-wide residue-site table must not infer maximal ideals, residue
fields, or tensor factors without certified local decomposition evidence. &
\texttt{Question} \\
\addlinespace
QSA Step 20: infinite boundary cases &
AQM-29, AQM-30, AQM-45, and AQM-49 separate finite supports, quasi-local
algebras, regular-function profiles, and generic-point caveats. &
Infinite-volume field operators, generic-point sectors, local-field or
adelic completions, and \(C^*\)- or Hilbert-space completions still require
separate conventions and sources. &
\texttt{Gap} \\
\bottomrule
\end{tabular}
\end{center}
